\documentclass{szb-solution}

\area{Analitikus geometria}
\title{Vektorok}
\subject{Matematika G1}
\subjectCode{BMETE94BG01}
\date{Utoljára frissítve: \today}
\docno{1}

\begin{document}
\maketitle

\subsection{Kollinearitás paraméterekkel}

Legyen $\rvec u$ és $\rvec v$ két tetszőleges vektor. Milyen $\alpha$ és
$\beta$ paraméterek esetén lesznek kollineárisak, ha az
$\{\rvec a;\rvec b;\rvec c\}$ vektorrendszer lineárisan független?

\begin{enumerate}[a)]
  \item
        $\rvec u=2\rvec a+3\rvec b$ és
        $\rvec v=4\rvec a+\alpha\rvec b$.

        A kollinearitás feltétele, hogy valamely $\lambda\in\Reals$ esetén
        $\rvec u=\lambda\rvec v$. A lineáris függetlenség miatt az együtthatók
        rendre megegyeznek:
        $$
          2=4\lambda,
          \qquad
          3=\alpha\lambda
          \quad\Longrightarrow\quad
          \lambda=\frac12,
          \qquad
          \boxed{\alpha=6}
          \text.
        $$

  \item
        $\rvec u=3\rvec a-3\alpha\rvec b+\beta\rvec c$ és
        $\rvec v=\rvec a-\alpha\rvec b-\rvec c$.

        Ismét $\rvec u=\lambda\rvec v$ alakban hasonlítjuk össze az
        együtthatókat:
        $$
          3=\lambda,
          \qquad
          -3\alpha=-\lambda\alpha,
          \qquad
          \beta=-\lambda
          \text.
        $$
        Innen $\lambda=3$, a középső egyenlet minden $\alpha$ esetén teljesül,
        ezért
        $$
          \boxed{\alpha\in\Reals,
          \qquad
          \beta=-3}
          \text.
        $$
\end{enumerate}

\subsection{Lineáris függetlenség}

Legyen a $\{\rvec a;\rvec b;\rvec c\}$ vektorrendszer lineárisan független.
Lineárisan független lesz-e a $\{\rvec r;\rvec s;\rvec t\}$ vektorrendszer?

\begin{enumerate}[a)]
  \item Ha
        $$
          \rvec r=3\rvec a+2\rvec b+\rvec c,
          \qquad
          \rvec s=5\rvec a-3\rvec b-2\rvec c,
          \qquad
          \rvec t=\nvec,
        $$
        akkor a rendszer a nullvektor miatt
        \textbf{lineárisan összefüggő}.

  \item Ha
        $$
          \rvec r=\rvec a+\rvec b+\rvec c,
          \qquad
          \rvec s=\rvec b+\rvec c,
          \qquad
          \rvec t=\rvec a+\rvec c,
        $$
        akkor a
        $\lambda_1\rvec r+\lambda_2\rvec s+\lambda_3\rvec t=\nvec$
        egyenletből
        $$
          (\lambda_1+\lambda_3)\rvec a
          +(\lambda_1+\lambda_2)\rvec b
          +(\lambda_1+\lambda_2+\lambda_3)\rvec c=\nvec
        $$
        adódik. Az eredeti rendszer lineáris függetlensége miatt
        $$
          \lambda_1+\lambda_3=0,
          \qquad
          \lambda_1+\lambda_2=0,
          \qquad
          \lambda_1+\lambda_2+\lambda_3=0
          \text.
        $$
        Ennek egyetlen megoldása
        $\lambda_1=\lambda_2=\lambda_3=0$, tehát a rendszer
        \textbf{lineárisan független}.
\end{enumerate}

\subsection{Vektorok végpontjainak kollinearitása}

Legyen $\rvec a$, $\rvec b$ és $\rvec c$ három közös kezdőpontú komplanáris
vektor, ahol $\rvec a$ és $\rvec b$ nem kollineáris. Bizonyítsuk be, hogy a
végpontjaik akkor és csak akkor esnek egy egyenesre, ha
$\rvec c=\alpha\rvec a+\beta\rvec b$ előállításban $\alpha+\beta=1$.

Ha a három végpont egy egyenesen van, akkor valamely $t\in\Reals$ esetén
$$
  \rvec c=\rvec a+t(\rvec b-\rvec a)
  =(1-t)\rvec a+t\rvec b
  \text.
$$
Ezért $\alpha=1-t$, $\beta=t$, vagyis $\alpha+\beta=1$.

Megfordítva, ha $\alpha+\beta=1$, akkor $\alpha=1-\beta$, így
$$
  \rvec c=(1-\beta)\rvec a+\beta\rvec b
  =\rvec a+\beta(\rvec b-\rvec a)
  \text.
$$
Ez éppen az $\rvec a$ és $\rvec b$ végpontján átmenő egyenes paraméteres
egyenlete, tehát $\rvec c$ végpontja is ezen az egyenesen van.

\subsection{Két vektor által bezárt szög}

Legyen $\rvec a=(7;-1;6)$ és $\rvec b=(2;20;2)$. Ekkor
$$
  \rvec a\cdot\rvec b=7\cdot2+(-1)\cdot20+6\cdot2=6,
$$
$$
  |\rvec a|=\sqrt{7^2+(-1)^2+6^2}=\sqrt{86},
  \qquad
  |\rvec b|=\sqrt{2^2+20^2+2^2}=\sqrt{408}
  \text.
$$
Ezért
$$
  \cos\varphi
  =\frac{\rvec a\cdot\rvec b}{|\rvec a|\,|\rvec b|}
  =\frac{6}{\sqrt{86}\sqrt{408}},
  \qquad
  \boxed{\varphi\approx88{,}16^\circ}
  \text.
$$

\subsection{Merőlegességi feltétel}

A $\rvec b=(6;-2;z)$ vektor akkor merőleges az
$\rvec a=(2;-3;1)$ vektorra, ha skaláris szorzatuk nulla:
$$
  \rvec a\cdot\rvec b
  =2\cdot6+(-3)(-2)+z
  =18+z=0
  \quad\Longrightarrow\quad
  \boxed{z=-18}
  \text.
$$

\subsection{A bezárt szög koszinusza}

A feltételekből
\begin{align*}
  (\rvec a+3\rvec b)\cdot(7\rvec a-5\rvec b)=0
  &\quad\Longrightarrow\quad
  7|\rvec a|^2+16\rvec a\cdot\rvec b-15|\rvec b|^2=0,
  \\
  (\rvec a-4\rvec b)\cdot(7\rvec a-2\rvec b)=0
  &\quad\Longrightarrow\quad
  7|\rvec a|^2-30\rvec a\cdot\rvec b+8|\rvec b|^2=0
  \text.
\end{align*}
Tehát
$$
  \rvec a\cdot\rvec b
  =\frac{-7|\rvec a|^2+15|\rvec b|^2}{16}
  =\frac{7|\rvec a|^2+8|\rvec b|^2}{30}
  \text,
$$
amiből $|\rvec a|^2=|\rvec b|^2$ következik. Ezt visszahelyettesítve
$\rvec a\cdot\rvec b=|\rvec a|^2/2$, ezért
$$
  \boxed{
    \cos\varphi
    =\frac{\rvec a\cdot\rvec b}{|\rvec a|\,|\rvec b|}
    =\frac12
  }
  \text.
$$

\subsection{A téglalap hiányzó csúcsa}

Az ismert csúcsok $A(2;6;0)$, $B(1;2;3)$ és $C(-2;8;z)$. A $B$ csúcsból
induló oldalirányok:
$$
  \overrightarrow{BA}=A-B=(1;4;-3),
  \qquad
  \overrightarrow{BC}=C-B=(-3;6;z-3)
  \text.
$$
A téglalap oldalai merőlegesek, ezért
$$
  \overrightarrow{BA}\cdot\overrightarrow{BC}
  =-3+24-3(z-3)=30-3z=0
  \quad\Longrightarrow\quad
  \boxed{z=10}
  \text.
$$
Ekkor $\overrightarrow{BC}=(-3;6;7)$ és
$\overrightarrow{AD}=\overrightarrow{BC}$, tehát
$$
  D=A+\overrightarrow{BC}
  =(2;6;0)+(-3;6;7)
  =\boxed{(-1;12;7)}
  \text.
$$

\subsection{Keresztszorzat}

Az $\rvec a=(-4;2;1)$ és $\rvec b=(-2;7;8)$ vektorok keresztszorzata:
$$
  \rvec a\times\rvec b
  =\begin{bmatrix}
    2\cdot8-1\cdot7 \\
    1\cdot(-2)-(-4)\cdot8 \\
    (-4)\cdot7-2\cdot(-2)
  \end{bmatrix}
  =\boxed{
    \begin{bmatrix}
      9\\30\\-24
    \end{bmatrix}
  }
  \text.
$$

\subsection{Keresztszorzatos kifejezés egyszerűsítése}

A bilinearitást és az antikommutativitást használva:
\begin{align*}
  (3\rvec a-\rvec b)\times(\rvec b+3\rvec a)
  &=3(\rvec a\times\rvec b)
    +9(\rvec a\times\rvec a)
    -(\rvec b\times\rvec b)
    -3(\rvec b\times\rvec a)
  \\
  &=3(\rvec a\times\rvec b)+3(\rvec a\times\rvec b)
  =\boxed{6(\rvec a\times\rvec b)}
  \text.
\end{align*}

\subsection{Kollinearitás koordinátákból}

Az $\rvec a=(-3;4;7)$ és $\rvec b=(2;5;1)$ vektorokra
$$
  \rvec a\times\rvec b
  =\begin{bmatrix}
    4\cdot1-7\cdot5 \\
    7\cdot2-(-3)\cdot1 \\
    (-3)\cdot5-4\cdot2
  \end{bmatrix}
  =\begin{bmatrix}
    -31\\17\\-23
  \end{bmatrix}
  \neq\nvec
  \text.
$$
Így a két vektor \textbf{nem kollineáris}.

\subsection{Háromszög területe}

Az $A(1;0;2)$, $B(-1;4;7)$ és $C(5;-2;1)$ pontok esetén
$$
  \overrightarrow{AB}=(-2;4;5),
  \qquad
  \overrightarrow{AC}=(4;-2;-1)
  \text,
$$
$$
  \overrightarrow{AB}\times\overrightarrow{AC}
  =\begin{bmatrix}
    6\\18\\-12
  \end{bmatrix},
  \qquad
  \left|\overrightarrow{AB}\times\overrightarrow{AC}\right|
  =\sqrt{6^2+18^2+(-12)^2}=6\sqrt{14}
  \text.
$$
A háromszög területe a kifeszített paralelogramma területének fele:
$$
  \boxed{T=\frac12\cdot6\sqrt{14}=3\sqrt{14}}
  \text.
$$

\subsection{Következik-e a vektorok egyenlősége?}

Az
$\rvec a\times\rvec c=\rvec b\times\rvec c$ egyenletből
$$
  (\rvec a-\rvec b)\times\rvec c=\nvec
$$
következik. Ez csak azt jelenti, hogy $\rvec a-\rvec b$ és $\rvec c$
kollineárisak (vagy $\rvec c=\nvec$), tehát általában
\textbf{nem következik}, hogy $\rvec a=\rvec b$.

Például legyen $\rvec c=(1;0;0)$, $\rvec a=(1;0;0)$ és
$\rvec b=(0;0;0)$. Ekkor mindkét keresztszorzat nulla, miközben
$\rvec a\neq\rvec b$.

\subsection{Egy kocka élvektorai}

Legyen $\rvec a=(6;2;-3)$ és $\rvec b=(-3;6;-2)$. Ekkor
$$
  \rvec a\cdot\rvec b
  =6(-3)+2\cdot6+(-3)(-2)=0,
$$
$$
  |\rvec a|=\sqrt{36+4+9}=7,
  \qquad
  |\rvec b|=\sqrt{9+36+4}=7
  \text.
$$
A két vektor tehát merőleges és egyenlő hosszú, így lehetnek egy kocka
élvektorai. Továbbá
$$
  \rvec a\times\rvec b
  =\begin{bmatrix}14\\21\\42\end{bmatrix}
  =7\begin{bmatrix}2\\3\\6\end{bmatrix}
  \text,
$$
ahol $(2;3;6)$ hossza $7$. A harmadik élvektor ezért
$$
  \boxed{\rvec c=\pm(2;3;6)}
  \text.
$$

\subsection{Három vektor lineáris függetlensége}

Legyen $\rvec a=(2;3;-1)$, $\rvec b=(1;-1;3)$ és
$\rvec c=(1;9;-11)$. A vegyesszorzat:
$$
  \rvec b\times\rvec c
  =\begin{bmatrix}-16\\14\\10\end{bmatrix},
  \qquad
  \rvec a\cdot(\rvec b\times\rvec c)
  =2(-16)+3\cdot14+(-1)\cdot10=0
  \text.
$$
A vegyesszorzat nulla, tehát a három vektor
\textbf{lineárisan összefüggő}.

\subsection{Paralelepipedon térfogata}

Az $\rvec a$, $\rvec b$ és $\rvec c$ által kifeszített paralelepipedon
térfogata legyen $V$. Az új élvektorok együtthatómátrixa
$$
  \rmat M=
  \begin{bmatrix}
    2&1&2\\
    3&-1&4\\
    4&1&-1
  \end{bmatrix},
  \qquad
  \det\rmat M=27
  \text,
$$
mivel az oszlopok rendre
$\rvec r=2\rvec a+3\rvec b+4\rvec c$,
$\rvec s=\rvec a-\rvec b+\rvec c$ és
$\rvec t=2\rvec a+4\rvec b-\rvec c$ koordinátái.
Az új térfogat ezért
$$
  \boxed{V_{\rvec r,\rvec s,\rvec t}
  =|\det\rmat M|\,V=27V}
  \text.
$$

\subsection{Paraméteres lineáris függetlenség}

Legyen $\rvec a=(3;\alpha;0)$, $\rvec b=(0;3;\alpha)$ és
$\rvec c=(1;0;-1)$. Ekkor
$$
  \rvec b\times\rvec c
  =\begin{bmatrix}-3\\\alpha\\-3\end{bmatrix},
  \qquad
  \rvec a\cdot(\rvec b\times\rvec c)
  =-9+\alpha^2
  \text.
$$
A rendszer pontosan akkor lineárisan összefüggő, ha a vegyesszorzat nulla:
$$
  \boxed{\alpha=\pm3
  \quad\Longrightarrow\quad
  \text{lineárisan összefüggő}}
  \text.
$$
Minden más valós paraméterérték esetén
$$
  \boxed{\alpha\in\Reals\setminus\{-3,3\}
  \quad\Longrightarrow\quad
  \text{lineárisan független}}
  \text.
$$

\subsection{Vektor vetülete}

Legyen $\rvec a=(-1;2;1)$ és $\rvec b=(1;2;2)$. Mivel
$$
  \rvec a\cdot\rvec b=-1+4+2=5,
  \qquad
  |\rvec b|^2=1+4+4=9,
$$
az $\rvec a$ vektor $\rvec b$ irányú vetülete
$$
  \boxed{
    \operatorname{proj}_{\rvec b}\rvec a
    =\frac{\rvec a\cdot\rvec b}{|\rvec b|^2}\rvec b
    =\frac59\begin{bmatrix}1\\2\\2\end{bmatrix}
    =\begin{bmatrix}5/9\\10/9\\10/9\end{bmatrix}
  }
  \text.
$$
A merőleges komponens:
$$
  \rvec a_\perp
  =\rvec a-\operatorname{proj}_{\rvec b}\rvec a
  =\begin{bmatrix}-14/9\\8/9\\-1/9\end{bmatrix}
  \text.
$$

\end{document}
